% Source-aligned English for printed p.197.
\clearpage
\SourcePage{197}

\LayerTarget{tohoku:conclusion:injective-g-o-embedding:p1}
purpose of obtaining a convenient explicit form for ``enough'' injective
objects, by generalizing the construction made in 3.1.
\LayerTarget{tohoku:construction:pointwise-g-o-action}
Let $(A_x)_{x\in X}$ be a family in which $A_x$ is an ${\mathbf O}(x)$-module, and
consider the product ${\mathbf O}$-module that it defines (cf. 3.1). To make $G$ act on
it so as to obtain a $G$-${\mathbf O}$-sheaf, it is enough, for every $x\in X$ and
$g\in G$, to give a homomorphism $a\to g.\,a$ from the
abelian group $A_x$ to the abelian group $A_{g.\, x}$ such that $g.\,(ua)=(g.\,u)(g.\,a)$ for $u\in {\mathbf O}(x)$ and $a\in A_x$, and $e.\,a=a$, $g.\,(g'.\,a)=(gg').\,a$.
\LayerTarget{tohoku:construction:auxiliary-ring-u-x}
Now introduce the auxiliary ring $U_x$ generated by the algebra
$\mathbf Z(G_x)$ of $G_x$ and the ring ${\mathbf O}(x)$, subject to the commutation
relations $gug^{-1}=u^g$ for $u\in {\mathbf O}(x)$, $g\in G_x$,
where, to avoid confusion, $u^g$ denotes the transform of $u$ by $g\in G_x$.
Then the $A_x$ are in fact $U_x$-modules, and any $g\in G$ defines an
isomorphism $a\to g.\,a$ from the group $A_x$ onto
$A_{g.\, x}$ satisfying $g.\,(ua)=(g.\,u)(g.\,a)$ for $u\in U_x$, $a\in A_x$.
\LayerTarget{tohoku:construction:orbit-representative-u-y}
It follows readily that, if an element $\xi(y)$ is chosen in every orbit
$y\in Y$ and one sets $U_y=U_{\xi(y)}$, the preceding data are equivalent to
the data of a family $(A_y)_{y\in Y}$ of $U_y$-modules
$A_y$ $(=A_{\xi(y)})$.

\LayerTarget{tohoku:construction:induced-a-bar-y}
The $A_x$ for $x\in y$ can be recovered from $A_y$ as follows. Let
$\overline A_y$ be the $G$-module \textit{induced} by the $G_y$-module $A_y$,
i.e. the $G$-module $\mathrm{Hom}_{G_y}(\mathbf Z(G),A_y)$ obtained from the
$G_y$-module $A_y$ by ``contravariant extension of scalars''. Then $A_y$
identifies with a quotient group of $\overline A_y$ by an abelian subgroup
stable under $G_y$, while $\overline A_y$ identifies with the product of all
the transformed quotients $g.\,A_y$ for $g\in G/G_y$.
\LayerTarget{tohoku:construction:induced-o-bar-y}
Likewise introduce the $G$-module $\overline O_y$ induced by
${\mathbf O}_y={\mathbf O}(\xi(y))$. One sees at once that $\overline A_y$ is an
$\overline O_y$-module identifying with the product of the
$g.\,O_y$-modules $g.\,A_y$ for $g\in G/G_y$. There
are therefore canonical isomorphisms $g.\,{\mathbf O}_y={\mathbf O}(g.\,\xi(y))$, $g.\,A_y=A_{g.\,\xi(y)}$.
\LayerTarget{tohoku:construction:p-a-g-o-sheaf}
Let $P(A)$ denote the $G$-${\mathbf O}$-sheaf on $X$ defined by the family
$A=(A_y)$ of $U_y$-modules $A_y$.
\LayerTarget{tohoku:comparison:hom-to-p-a}
Let $B$ be any $G$-${\mathbf O}$-sheaf. An ${\mathbf O}$-homomorphism from a $G$-${\mathbf O}$-module $B$
to $P(A)$ identifies with a family $(v_x)_{x\in X}$ of
${\mathbf O}(x)$-homomorphisms $B(x)\to A_x$ (as noted in 3.1), and this
${\mathbf O}$-homomorphism is compatible with $G$ if and only if the $v_x$ ``commute''
with $G$. It is then immediate that giving a $G$-${\mathbf O}$-homomorphism from $B$ to
$P(A)$ is equivalent to giving a family $(v_y)_{y\in Y}$ of
$U_y$-homomorphisms $B(\xi(y))\to A_y$. We conclude at once:

\LayerTarget{tohoku:proposition:5.1.2}
\textsc{Proposition 5.1.2.} \textit{Let $(A_y)_{y\in Y}$ be a
family of injective $U_y$-modules. Then the $G$-${\mathbf O}$-module $P((A_y))$ that it
defines is injective. Moreover, every $G$-${\mathbf O}$-module is isomorphic to a
subsheaf of a sheaf of the preceding type.}

\LayerTarget{tohoku:proof:5.1.2}
The latter assertion, which is the more important one for us, is immediate:
for every $y$, embed the stalk module $B(\xi(y))$ of the given sheaf $B$ into
an injective $U_y$-module.

\LayerTarget{tohoku:construction:direct-image-g-sheaf:p1}
\textit{Direct images of $G$-sheaves.} Let $A$ be a sheaf of sets on
$X$. Recall that its direct image $f_*(A)$ was defined as the sheaf on $Y$
whose sections over an open subset $U\subset Y$ are the sections of $A$ over
$f^{-1}(U)$. If $A$ is a $G$-sheaf, then $\Gamma(f^{-1}(U),A)$ has $G$ as an
operator group, whence it follows readily that \textit{the direct image
$f_*(A)$ of a $G$-sheaf on $X$ is a $G$-sheaf on $Y$} (recall that $G$ acts
trivially on $Y$); of

\TohokuContinuation{198}
